% GENERATED FILE. Edit canonical lessons, then run npm run generate:books.
% canonical-source-hash: fnv1a64:fac220b0cb1177b9
% canonical-lessons: PA-C186-future-return, PA-C186-hovega, PA-C186-hovegi, PA-C186-hovega-hovegi-review, PA-C186-aavega, PA-C186-aavegi, PA-C186-future-review, PA-C186-payoff

\chapter{What Will Happen? Who Will Come?}
\label{ch:pa-a2-third-person-future}
\hlchaptermodality{186}

\begin{chapteropening}
\textbf{By the end of this chapter:} \emph{I can recognise and use common masculine and feminine third-person future forms of to be and to come.}

Retrieve four exact third-person future forms through listening, speaking, reading, and bounded copying after a slow agreement ramp.
\end{chapteropening}

\section[(future base review)]{Start with two verbs you know}

\label{lesson:PA-C186-future-return}



\begin{quote}
Say the Punjabi infinitive for “to be,” then the infinitive for “to come.”
\end{quote}

\begin{grammarlens}[title={Grammar Lens: one future contrast at a time}]
Keep the two familiar meanings steady. The next steps will make a masculine future form and then a feminine future form from the first verb. Only after that contrast is secure will the same two endings move to the second verb.
\end{grammarlens}

\subsection*{Your turn}
\begin{itemize}
\raggedright
  \item \emph{Hear:} \emph{honaa}; YOU SAY: \textbf{\pa{ਹੋਣਾ}}
  \item \emph{Hear:} \emph{aaunaa}; YOU SAY: \textbf{\pa{ਆਉਣਾ}}
\end{itemize}

\begin{tcolorbox}[breakable,colback=teal!4,colframe=teal!35!black,title={Before you move on}]
Say “to be; to come” once more in Punjabi.

This voice-first review introduces no new written target.
\end{tcolorbox}
\section[\texorpdfstring{\pa{ਹੋਵੇਗਾ} (hovegaa)}{ਹੋਵੇਗਾ (hovegaa)}]{\pa{ਹੋਵੇਗਾ} — it will be or happen}

\label{lesson:PA-C186-hovega}



\begin{quote}
Say \textbf{\pa{ਹੋਣਾ}}, then say \textbf{\pa{ਕੰਮ}}.
\end{quote}

\begin{grammarlens}[title={Grammar Lens: the masculine singular future}]
\begin{quote}
\textbf{\pa{ਹੋਵੇਗਾ}} — \emph{hovegaa} — it will be or it will happen
\end{quote}

This is the masculine singular future form. First keep it whole; the feminine contrast comes in the next small step.
\end{grammarlens}

\subsection*{Your turn}
\begin{itemize}
\raggedright
  \item \emph{Hear:} \emph{hovegaa}; YOU SAY: \textbf{\pa{ਹੋਵੇਗਾ}}
  \item \emph{Read it:} \textbf{\pa{ਕੰਮ} \pa{ਹੋਵੇਗਾ}}
  \item \emph{Write it:} copy \textbf{\pa{ਹੋਵੇਗਾ}} once and check the final \textbf{\pa{ਾ}}
\end{itemize}

\begin{tcolorbox}[breakable,colback=teal!4,colframe=teal!35!black,title={Before you move on}]
Read the exact form for “it will happen.” (\textbf{\pa{ਹੋਵੇਗਾ}}.)

Source: standard Punjabi uses this form as the masculine singular future of the already-taught infinitive \textbf{\pa{ਹੋਣਾ}}.
\end{tcolorbox}
\section[\texorpdfstring{\pa{ਹੋਵੇਗੀ} (hovegii)}{ਹੋਵੇਗੀ (hovegii)}]{\pa{ਹੋਵੇਗੀ} — it will be, with a feminine subject}

\label{lesson:PA-C186-hovegi}



\begin{quote}
Say \textbf{\pa{ਹੋਵੇਗਾ}}, then say \textbf{\pa{ਯਾਤਰਾ}}.
\end{quote}

\begin{grammarlens}[title={Grammar Lens: change only the agreement ending}]
\begin{quote}
\textbf{\pa{ਹੋਵੇਗੀ}} — \emph{hovegii} — it will be or happen, feminine singular
\end{quote}

The ending changes from \textbf{\pa{ਾ}} to \textbf{\pa{ੀ}}. The time stays future; the form now agrees with a feminine singular subject.
\end{grammarlens}

\subsection*{Your turn}
\begin{itemize}
\raggedright
  \item \emph{Hear:} \emph{hovegaa, hovegii}; YOU SAY: \textbf{\pa{ਹੋਵੇਗਾ}, \pa{ਹੋਵੇਗੀ}}
  \item \emph{Read it:} \textbf{\pa{ਯਾਤਰਾ} \pa{ਹੋਵੇਗੀ}}
  \item \emph{Write it:} copy \textbf{\pa{ਹੋਵੇਗੀ}} once and check the final \textbf{\pa{ੀ}}
\end{itemize}

\begin{tcolorbox}[breakable,colback=teal!4,colframe=teal!35!black,title={Before you move on}]
Say the masculine form, then the feminine form.

Source: standard Punjabi future agreement changes the final masculine \textbf{\pa{ਾ}} to feminine \textbf{\pa{ੀ}} here.
\end{tcolorbox}
\section[\texorpdfstring{\pa{ਹੋਵੇਗਾ} · \pa{ਹੋਵੇਗੀ} (hovegaa / hovegii)}{ਹੋਵੇਗਾ · ਹੋਵੇਗੀ (hovegaa / hovegii)}]{Hold the first future contrast steady}

\label{lesson:PA-C186-hovega-hovegi-review}



\begin{quote}
Say the pair once without looking.
\end{quote}

\subsection*{Your turn}
\begin{itemize}
\raggedright
  \item \emph{Hear:} \emph{kamm hovegaa}; YOU SAY: \textbf{\pa{ਕੰਮ} \pa{ਹੋਵੇਗਾ}}
  \item \emph{Hear:} \emph{yaatraa hovegii}; YOU SAY: \textbf{\pa{ਯਾਤਰਾ} \pa{ਹੋਵੇਗੀ}}
  \item \emph{Read it:} \textbf{\pa{ਹੋਵੇਗਾ} · \pa{ਹੋਵੇਗੀ}}
  \item \emph{Write it:} cover the pair and write \textbf{\pa{ਹੋਵੇਗਾ}, \pa{ਹੋਵੇਗੀ}}
\end{itemize}

The stem and future meaning stay together. Only the final agreement vowel changes.

\begin{tcolorbox}[breakable,colback=teal!4,colframe=teal!35!black,title={Before you move on}]
Which form ends in \textbf{\pa{ਾ}}? Which ends in \textbf{\pa{ੀ}}?

This review adds no new future form.
\end{tcolorbox}
\section[\texorpdfstring{\pa{ਆਵੇਗਾ} (aavegaa)}{ਆਵੇਗਾ (aavegaa)}]{\pa{ਆਵੇਗਾ} — he or it will come}

\label{lesson:PA-C186-aavega}



\begin{quote}
Say \textbf{\pa{ਆਉਣਾ}}, then retrieve the masculine future ending from the last pair.
\end{quote}

\begin{grammarlens}[title={Grammar Lens: transfer the masculine future}]
\begin{quote}
\textbf{\pa{ਆਵੇਗਾ}} — \emph{aavegaa} — he or it will come
\end{quote}

The verb changes, but the masculine singular future ending remains familiar.
\end{grammarlens}

\subsection*{Your turn}
\begin{itemize}
\raggedright
  \item \emph{Hear:} \emph{aavegaa}; YOU SAY: \textbf{\pa{ਆਵੇਗਾ}}
  \item \emph{Read it:} \textbf{\pa{ਅਮਨ} \pa{ਆਵੇਗਾ}}
  \item \emph{Write it:} copy \textbf{\pa{ਆਵੇਗਾ}} once and check the final \textbf{\pa{ਾ}}
\end{itemize}

\begin{tcolorbox}[breakable,colback=teal!4,colframe=teal!35!black,title={Before you move on}]
Read “will come,” masculine singular. (\textbf{\pa{ਆਵੇਗਾ}}.)

Source: standard Punjabi uses this form as the masculine singular future of the already-taught infinitive \textbf{\pa{ਆਉਣਾ}}.
\end{tcolorbox}
\section[\texorpdfstring{\pa{ਆਵੇਗੀ} (aavegii)}{ਆਵੇਗੀ (aavegii)}]{\pa{ਆਵੇਗੀ} — she or it will come}

\label{lesson:PA-C186-aavegi}



\begin{quote}
Say \textbf{\pa{ਆਵੇਗਾ}}, then say \textbf{\pa{ਰੇਲ}}.
\end{quote}

\begin{grammarlens}[title={Grammar Lens: transfer the feminine future}]
\begin{quote}
\textbf{\pa{ਆਵੇਗੀ}} — \emph{aavegii} — she or it will come
\end{quote}

Change the final \textbf{\pa{ਾ}} to \textbf{\pa{ੀ}}, just as in the first future pair. The verb still means “come”; agreement is now feminine singular.
\end{grammarlens}

\subsection*{Your turn}
\begin{itemize}
\raggedright
  \item \emph{Hear:} \emph{aavegaa, aavegii}; YOU SAY: \textbf{\pa{ਆਵੇਗਾ}, \pa{ਆਵੇਗੀ}}
  \item \emph{Read it:} \textbf{\pa{ਰੇਲ} \pa{ਆਵੇਗੀ}}
  \item \emph{Write it:} copy \textbf{\pa{ਆਵੇਗੀ}} once and check the final \textbf{\pa{ੀ}}
\end{itemize}

\begin{tcolorbox}[breakable,colback=teal!4,colframe=teal!35!black,title={Before you move on}]
Say the masculine “will come,” then the feminine “will come.”

Source: standard Punjabi future agreement changes the final masculine \textbf{\pa{ਾ}} to feminine \textbf{\pa{ੀ}} here too.
\end{tcolorbox}
\section[\texorpdfstring{\pa{ਹੋਵੇਗਾ} · \pa{ਹੋਵੇਗੀ} … (hovegaa …)}{ਹੋਵੇਗਾ · ਹੋਵੇਗੀ … (hovegaa …)}]{Two verbs, one future agreement contrast}

\label{lesson:PA-C186-future-review}



\begin{quote}
Say both masculine forms, then both feminine forms.
\end{quote}

\subsection*{Your turn}
\begin{itemize}
\raggedright
  \item \emph{Hear:} \emph{kamm hovegaa; yaatraa hovegii}; YOU SAY: \textbf{\pa{ਕੰਮ} \pa{ਹੋਵੇਗਾ}; \pa{ਯਾਤਰਾ} \pa{ਹੋਵੇਗੀ}}
  \item \emph{Hear:} \emph{Aman aavegaa; rel aavegii}; YOU SAY: \textbf{\pa{ਅਮਨ} \pa{ਆਵੇਗਾ}; \pa{ਰੇਲ} \pa{ਆਵੇਗੀ}}
  \item \emph{Read it:} \textbf{\pa{ਹੋਵੇਗਾ}, \pa{ਹੋਵੇਗੀ}, \pa{ਆਵੇਗਾ}, \pa{ਆਵੇਗੀ}}
  \item \emph{Write it:} cover the line and write the four forms in the same order
\end{itemize}

Keep the verb pair separate from the agreement pair: “be” and “come” each have one masculine and one feminine singular future form here.

\begin{tcolorbox}[breakable,colback=teal!4,colframe=teal!35!black,title={Before you move on}]
Retrieve all four exact forms without looking.

This review adds no first-person future and no past form.
\end{tcolorbox}
\section[\texorpdfstring{\pa{ਕੰਮ} \pa{ਹੋਵੇਗਾ} · \pa{ਯਾਤਰਾ} … (kamm hovegaa …)}{ਕੰਮ ਹੋਵੇਗਾ · ਯਾਤਰਾ … (kamm hovegaa …)}]{A gentle third-person future payoff}

\label{lesson:PA-C186-payoff}



\begin{quote}
Without looking, say one future form ending in \textbf{\pa{ਾ}} and one ending in \textbf{\pa{ੀ}} for each verb.
\end{quote}

\subsection*{Your turn: four skills, bounded to the two pairs}
\begin{itemize}
\raggedright
  \item \emph{Hear:} \emph{kamm hovegaa, yaatraa hovegii, Aman aavegaa, rel aavegii}; YOU SAY: \textbf{\pa{ਕੰਮ} \pa{ਹੋਵੇਗਾ}, \pa{ਯਾਤਰਾ} \pa{ਹੋਵੇਗੀ}, \pa{ਅਮਨ} \pa{ਆਵੇਗਾ}, \pa{ਰੇਲ} \pa{ਆਵੇਗੀ}}
  \item \emph{Read it:} \textbf{\pa{ਕੰਮ} \pa{ਹੋਵੇਗਾ।} \pa{ਯਾਤਰਾ} \pa{ਹੋਵੇਗੀ।} \pa{ਅਮਨ} \pa{ਆਵੇਗਾ।} \pa{ਰੇਲ} \pa{ਆਵੇਗੀ।}}
  \item \emph{Write it:} cover the model and write \textbf{\pa{ਹੋਵੇਗਾ}, \pa{ਹੋਵੇਗੀ}, \pa{ਆਵੇਗਾ}, \pa{ਆਵੇਗੀ}}
\end{itemize}

Now underline the two masculine endings and circle the two feminine endings.

\begin{tcolorbox}[breakable,colback=teal!4,colframe=teal!35!black,title={Before you move on}]
Retrieve the four measured forms once more.

This is the ninth prerequisite tranche exposed by Punjabi A2 mock drafting. It does not unblock Punjabi A2 mock 1 by itself. It is not an A2 mock, not A2 readiness or attainment, and not external or human validation.

Source: the standard Punjabi third-person future agreement pattern applied to the already-taught infinitives \textbf{\pa{ਹੋਣਾ}} and \textbf{\pa{ਆਉਣਾ}}.
\end{tcolorbox}
